% ============================================================
\section{Performance Comparison}
% ============================================================

\subsection{Latency Across Configurations}

Unless stated otherwise, all latency results in this section were obtained using \textbf{\texttt{cosf}} (single-precision) and compiler optimisation level \textbf{-O2} (with debug enabled), averaged over 20 runs. 

\begin{table}[H]
\centering
\caption{Average Latency (Ticks, -O2)}
\begin{tabular}{lcccc}
\toprule
Variant & Case1 & Case2 & Case3 & Case4 \\
\midrule
V3 (SW baseline) & 10    & 435   & 14381 & 505 \\
V4 (NiosII HW mul) & 5     & 204   & 6563  & 232 \\
V6 (FP hardware) & \multicolumn{4}{c}{\textit{[to be filled]}} \\
\bottomrule
\end{tabular}
\end{table}

The results show approximately linear scaling with $N$ across configurations, consistent with the $\mathcal{O}(N)$ structure of the summation. 
Enabling Nios~II hardware multiplier support (V4) provides a consistent $\sim 2\times$ speedup over the software baseline (V3), indicating that multiplication dominates a substantial portion of the runtime at this stage.

\subsection{Effect of \texttt{cos} vs \texttt{cosf}}

In this experiment, the data were collected under configuration \textbf{V4} (i.e., only Nios~II hardware multiplier support enabled, 3$\times$16-bit multipliers), compiled with \textbf{-O2}. 
All other computation remained software-based.

\begin{table}[H]
\centering
\caption{\texttt{cosf} vs \texttt{cos} (Ticks, V4, -O2)}
\begin{tabular}{lccc}
\toprule
Case & cosf & cos & Speedup (cosf vs cos) \\
\midrule
1 & 5    & 7    & 1.40$\times$ \\
2 & 204  & 309  & 1.51$\times$ \\
3 & 6563 & 9942 & 1.51$\times$ \\
4 & 232  & 356  & 1.53$\times$ \\
\bottomrule
\end{tabular}
\end{table}

Using \texttt{cosf} consistently reduces latency across all test cases, delivering a speedup of approximately \textbf{$1.4\times$ to $1.5\times$} compared to \texttt{cos}. 
This is expected because \texttt{cos} typically performs double-precision computation and introduces additional conversion overhead, whereas \texttt{cosf} evaluates the trigonometric function directly in single precision. 
The results indicate that trigonometric evaluation remains a major contributor to runtime even with multiplier acceleration, motivating further hardware support for floating-point and transcendental operations in later tasks.